\chapter{Cushing's Syndrome}
\index{Cushing's Syndrome}
\label{sec:Cushing's Syndrome}

\section{Overview}
This condition affects many people worldwide and can range from mild to severe. Recognizing the condition early and managing it properly gives you the best chance for a good outcome. Understanding what causes this condition helps your doctor choose the right treatment for you. Learning about your condition and taking an active role in your care are key to managing it long term.

\begin{itemize}[noitemsep]
  \item Cushing's syndrome results from chronic exposure to excess cortisol, regardless of the cause, with exogenous glucocorticoid use (iatrogenic) being the most common cause overall. Endogenous causes are divided into ACTH-dependent (70--80\%): Cushing's disease (pituitary adenoma secreting ACTH, 70\%) and ectopic ACTH syndrome (small cell lung cancer, bronchial carcinoid)
  \item And ACTH-independent (20--30\%): adrenal adenoma, adrenal carcinoma, and bilateral adrenal abnormal increase in cells.
\end{itemize}

\section{Causes}
This condition develops from a combination of genetic and environmental factors that disrupt normal body processes. Several factors can work together to trigger or make the condition worse.

\section{Risk Factors}

\begin{itemize}[noitemsep]
  \item Genetic predisposition and family history
  \item Advancing age
  \item Lifestyle factors (diet, exercise, smoking, alcohol)
  \item Environmental exposures
  \item Underlying medical conditions
  \item Medication use
\end{itemize}

\section{Signs and Symptoms}
You may experience central (truncal) obesity, moon facies, dorsocervical fat pad (``buffalo hump''), supraclavicular fat pads, thin extremities, purple striae ($>$1 cm wide), easy bruising, proximal muscle weakness, high blood pressure, hyperglycemia/diabetes, osteoporosis, menstrual irregularity, hirsutism, and psychiatric disturbances.

\begin{itemize}[noitemsep]
  \item Pain or discomfort in the affected area
  \item Difficulty performing daily activities
  \item Changes in how the affected area looks or works
  \item General symptoms may include fatigue, fever, or weight changes
  \item Symptoms may develop slowly or appear suddenly
\end{itemize}

\section{Progression and Stages}
This condition can progress through different stages over time. Early stages often involve mild symptoms that you may not notice right away. As the condition progresses, symptoms become more noticeable and may affect your daily activities. Advanced stages may involve more serious symptoms and complications.

\section{Complications}

\begin{itemize}[noitemsep]
  \item The condition may worsen over time if not treated
  \item Increased risk of other infections or injuries
  \item Side effects from medications or other treatments
  \item Reduced quality of life
  \item Emotional and psychological effects
\end{itemize}

\section{Diagnosis}
Diagnosis typically involves confirmation of hypercortisolism (24-hour urinary free cortisol, late-night salivary cortisol, or 1 mg overnight dexamethasone suppression test) followed by determination of the cause (plasma ACTH, high-dose dexamethasone suppression test, CRH stimulation test, MRI pituitary, CT adrenals, inferior petrosal sinus sampling).

\begin{itemize}[noitemsep]
  \item Your doctor will take a detailed medical history and perform a physical exam
  \item Lab tests to check how your organs are working
  \item Imaging tests (like X-rays or MRI) to look at the affected area
  \item Specialized tests depending on which body system is involved
  \item Referral to a specialist for further evaluation
\end{itemize}

\section{Prevention}

\begin{itemize}[noitemsep]
  \item Eat a balanced diet and exercise regularly
  \item Avoid known triggers and risk factors
  \item Go for regular check-ups so any problems are caught early
  \item Follow your doctor's screening recommendations
\end{itemize}

\section{Treatment Options}
The right treatment depends on the cause: transsphenoidal surgery for Cushing's disease, surgical removal of adrenal tumors, chemotherapy for adrenal carcinoma, and ketoconazole or metyrapone for medical management.

\begin{itemize}[noitemsep]
  \item Medications to control symptoms and treat the underlying cause
  \item Lifestyle changes and self-care
  \item Physical therapy and rehabilitation
  \item Surgery when needed
  \item Regular follow-up care
\end{itemize}

\section{Living with It}

\begin{itemize}[noitemsep]
  \item Follow your treatment plan as prescribed
  \item Keep up with regular appointments with your healthcare provider
  \item Adjust your daily habits as needed
  \item Reach out to family, friends, or support groups for help
  \item Keep learning about your condition and how to manage it
\end{itemize}

\section{Outlook}
The outlook is generally positive with successful treatment.
